\documentclass[11pt,a4paper]{article}

% ======================================================================
% CONFIGURACIÓN DEL PROYECTO Y PAQUETES
% ======================================================================
\usepackage[utf8]{inputenc}
\usepackage[T1]{fontenc}
\usepackage[spanish,es-tabla]{babel}
\usepackage[margin=2.5cm, headheight=30pt]{geometry} 
\usepackage{graphicx}
\usepackage{fancyhdr}
\usepackage{xcolor}
\usepackage[colorlinks=true, linkcolor=blue, urlcolor=blue, citecolor=black]{hyperref}

% --- Matemáticas y Electrónica ---
\usepackage{amsmath, amssymb}
\usepackage{siunitx}
\usepackage[american]{circuitikz}

% --- Elementos de Diseño y Código ---
\usepackage{tcolorbox}
\usepackage{booktabs}
\usepackage{enumitem}
\usepackage{listings}
\usepackage{tabularx}
\usepackage{microtype}

% ======================================================================
% MACROS Y ESTILOS PERSONALIZADOS
% ======================================================================
\sisetup{
    output-decimal-marker = {,},
    per-mode = symbol,
    inter-unit-product = \ensuremath{{}\cdot{}}
}

% --- Cajas Personalizadas ---
\definecolor{ucbblue}{RGB}{0,51,102}

\newtcolorbox{resultado}[1][]{
    colback=green!5!white,
    colframe=green!60!black,
    fonttitle=\bfseries,
    title=¿Qué deben esperar en esta Fase?,
    #1
}

\newtcolorbox{info}[1][]{
    colback=blue!5!white,
    colframe=ucbblue,
    fonttitle=\bfseries,
    title=Nota Técnica,
    #1
}

% --- Formato de Texto SPICE ---
\newcommand{\tecla}[1]{\colorbox{gray!20}{\texttt{\footnotesize\textbf{#1}}}}
\newcommand{\spice}[1]{\texttt{\color{purple}#1}}

% --- Configuración de Listings para Python ---
\lstset{
    language=Python,
    basicstyle=\ttfamily\footnotesize,
    keywordstyle=\color{blue}\bfseries,
    stringstyle=\color{teal},
    commentstyle=\color{green!50!black}\itshape,
    backgroundcolor=\color{gray!5},
    frame=single,
    rulecolor=\color{gray!40},
    breaklines=true,
    showstringspaces=false,
    numbers=left,
    numberstyle=\tiny\color{gray},
    numbersep=5pt
}

% ======================================================================
% ESTILO DE PÁGINA
% ======================================================================
\pagestyle{fancy}
\fancyhf{}
\lhead{\textbf{Circuitos Electrónicos II}}
\rhead{Ingeniería Mecatrónica -- UCB Tarija}
\cfoot{Página \thepage}
\renewcommand{\headrulewidth}{0.4pt}

\begin{document}

\begin{center}
    \rule{\textwidth}{1pt}\\[0.3cm]
    {\Large \textbf{LABORATORIO 1: PUENTE DE WHEATSTONE EN CA}}\\[0.2cm]
    {\large \textbf{Manual de Ejecución y Validación (Modelo F.I.V.E. / CDIO)}}\\[0.2cm]
    \rule{\textwidth}{1pt}
\end{center}
\vspace{0.3cm}

Este es el flujo de ejecución paso a paso para el Laboratorio 1, estructurado bajo el modelo F.I.V.E. (Fundamentación, Implementación, Validación, Evidencia) y el marco CDIO. En cada etapa se detalla la acción requerida, el procedimiento técnico y el resultado numérico/gráfico exacto que el estudiante debe obtener.

% ======================================================================
% FASE 0: FUNDAMENTACIÓN
% ======================================================================
\section*{FASE 0: Fundamentación Teórica Previa (Cálculo Analítico)}
Antes de abrir el software o energizar el circuito, los estudiantes deben calcular a mano o en un script base los valores de referencia para $T = [20, 25, 30, 40, 50, 60]\ ^\circ\text{C}$.

\textbf{Paso 0.1 — Resistencia del Sensor (Steinhart-Hart reducida):}
\begin{equation}
    R_{NTC}(T) = R_0 \cdot e^{\beta \left( \frac{1}{T} - \frac{1}{T_0} \right)}
\end{equation}
(con $R_0 = \qty{10}{\kilo\ohm}$, $\beta = \qty{3950}{\kelvin}$, $T_0 = \qty{298.15}{\kelvin}$, $T$ en Kelvin).

\textbf{Paso 0.2 — Impedancia Compleja de la Rama 4:}
A la frecuencia de portadora $f_c = \qty{10}{\kilo\hertz}$ ($\omega_c = 2\pi \cdot 10^4\text{ rad/s}$):
\begin{equation}
    Z_{sensor}(j\omega) = R_{cable} + j\omega L_{cable} + \left( \frac{R_{NTC}}{1 + j\omega C_{para} R_{NTC}} \right)
\end{equation}

\textbf{Paso 0.3 — Salida Fasorial Diferencial:}
\begin{equation}
    \hat{V}_{out} = \hat{V}_{in} \left( \frac{Z_{sensor}}{R_3 + Z_{sensor}} - \frac{R_2}{R_1 + R_2} \right)
\end{equation}

\begin{resultado}
A \qty{25}{\celsius}: $R_{NTC} = \qty{10.0}{\kilo\ohm} \implies Z_{sensor} \approx 3124\ \Omega \angle -67.1^\circ$.\\
Tensión diferencial teórica: $\vert{}\hat{V}_{out}\vert{} = \mathbf{181.5\text{ mV}_{pico}}$ con un desfase de $\theta = \mathbf{-162.1^\circ}$.\\[0.2cm]
A \qty{60}{\celsius}: $R_{NTC} \approx \qty{3.02}{\kilo\ohm} \implies \vert{}\hat{V}_{out}\vert{} = \mathbf{420.2\text{ mV}_{pico}}$.
\end{resultado}

% ======================================================================
% FASE 1: SIMULACIÓN VIRTUAL
% ======================================================================
\section*{FASE 1: Simulación Virtual en LTSpice}
El estudiante valida el modelo matemático y anticipa el comportamiento de las señales en el dominio del tiempo y la frecuencia.

\begin{center}
    \begin{circuitikz}[scale=0.9, transform shape]
        % Fuente movida más a la izquierda para evitar solapamiento
        \draw (-2,10) to[sV, l=$V_1$] (-2,0);
        \draw (-2,10) -- (2.5,10) -- (7.5,10);
        \draw (-2,0) -- (2.5,0) -- (7.5,0);
        \draw (2.5,0) node[ground]{};
        
        \draw (2.5,10) to[R, l=$R_1$] (2.5,8) node[circle,fill,inner sep=1.5pt]{};
        \draw (2.5,8) to[R, l=$R_2$] (2.5,0);
        \node[left, align=center] at (2,8) {Nodo A};
        
        \draw (7.5,10) to[R, l=$R_3$] (7.5,8) node[circle,fill,inner sep=1.5pt]{};
        \draw (7.5,8) to[L, l=$L_{cable}$, a=\qty{15}{\micro\henry}] (7.5,6) to[R, l=$R_{cable}$, a=\qty{10}{\ohm}] (7.5,4);
        \draw (7.5,4) -- (6.5,4) to[vR, l_=$R_{NTC}$] (6.5,2) -- (7.5,2);
        \draw (7.5,4) -- (8.5,4) to[C, l=$C_{para}$, a=\qty{4.7}{\nano\farad}] (8.5,2) -- (7.5,2);
        \draw (7.5,2) -- (7.5,0);
        \node[right, align=center] at (8.5,8) {Nodo B};
        
        \draw (2.5,8) to[short, -o] (4.5,8) node[above]{$+$};
        \draw (7.5,8) to[short, -o] (5.5,8) node[above]{$-$};
        \node at (5, 8) {$V_{diff}$};
    \end{circuitikz}
\end{center}

\textbf{Paso 1.1 — Configuración del esquemático:}
\begin{itemize}[itemsep=1pt]
    \item Configurar la fuente $V_1$: \param{SINE(0 1 10k)} y AC Amplitude: \param{1}.
    \item Etiquetar nodos como A y B usando la herramienta Label Net (tecla \tecla{n}).
    \item En el valor de R\_NTC, escribir: \spice{\{R0 * exp(Beta * (1/(Tc + 273.15) - 1/T0))\}}.
\end{itemize}

\textbf{Paso 1.2 — Inserción de Directivas SPICE (Tecla \tecla{.}):}
\begin{lstlisting}[language=SPICE, numbers=none]
.param R0=10k Beta=3950 T0=298.15 Tc=25
.step param Tc 20 60 10
.tran 0.5m
\end{lstlisting}

\textbf{Paso 1.3 — Trazado y análisis:}
\begin{itemize}[itemsep=1pt]
    \item Hacer clic en el Nodo A, arrastrar la sonda y soltar en el Nodo B para graficar \param{V(A,B)}.
    \item Activar dos cursores sobre la traza transitoria para medir la diferencia de tiempo ($\Delta t$) entre el cruce por cero de $V_{in}$ y el de $V(A,B)$.
\end{itemize}

\begin{resultado}
Una familia de 5 ondas senoidales de \qty{10}{\kilo\hertz} superpuestas (correspondientes a los pasos de \qty{20}{\celsius} a \qty{60}{\celsius}).\\
La amplitud pico de $V(A,B)$ crece progresivamente desde $\approx \qty{150}{\milli\volt}$ hasta $\approx \qty{420}{\milli\volt}$.\\
Un retraso temporal medido con cursores de $\Delta t \approx \qty{45.0}{\micro\second}$, lo que confirma analíticamente:
\begin{equation}
    \theta = \Delta t \cdot f_c \cdot 360^\circ = 45\ \mu\text{s} \cdot 10\text{ kHz} \cdot 360^\circ = 162.0^\circ
\end{equation}
\end{resultado}

\newpage
% ======================================================================
% FASE 2: IMPLEMENTACIÓN FÍSICA
% ======================================================================
\section*{FASE 2: Implementación Física y Adquisición en Hardware}
El estudiante traslada el diseño al protoboard, conecta los instrumentos de banco y adquiere las trazas reales.

\textbf{Paso 2.1 — Caracterización previa con Multímetro (DMM):}
\begin{itemize}[itemsep=1pt]
    \item Medir con el DMM los valores reales de $R_1, R_2, R_3$ y $R_{NTC}$ a temperatura ambiente.
    \item Calcular la temperatura ambiente real del laboratorio $T_{amb}$ despejando la fórmula de Steinhart-Hart a partir de la resistencia medida.
\end{itemize}

\textbf{Paso 2.2 — Montaje y Alambrado:}
\begin{itemize}[itemsep=1pt]
    \item Armar el puente en el protoboard incluyendo en la rama 4 los componentes parásitos reales: $R_{cable} = \qty{10}{\ohm}$, $L_{cable} = \qty{15}{\micro\henry}$ y $C_{para} = \qty{4.7}{\nano\farad}$.
\end{itemize}

\textbf{Paso 2.3 — Configuración de Instrumentación:}
\begin{itemize}[itemsep=1pt]
    \item \textbf{Generador:} Onda senoidal, \qty{2.00}{\volt}$_{pp}$ (\qty{1.00}{\volt}$_{pico}$), frecuencia \qty{10.00}{\kilo\hertz}, offset \qty{0}{\volt}.
    \item \textbf{Osciloscopio:}
    \begin{itemize}[itemsep=1pt]
        \item Canal 1 (CH1): Conectado al Nodo A (punta $1\times$ o $10\times$ calibrada).
        \item Canal 2 (CH2): Conectado al Nodo B.
        \item Pinzas de tierra de ambas sondas: Conectadas únicamente a la tierra común (GND) del puente.
        \item Canal Matemático (MATH): Configurar la operación $M = CH1 - CH2$.
    \end{itemize}
\end{itemize}

\textbf{Paso 2.4 — Adquisición y Exportación:}
\begin{itemize}[itemsep=1pt]
    \item Registrar datos para al menos 4 condiciones térmicas (Ambiente, Contacto con dedos $\approx \qty{35}{\celsius}$, Fuente de calor suave $\approx \qty{45}{\celsius}$ y $\approx \qty{55}{\celsius}$).
    \item Para cada punto térmico: insertar la memoria USB en el osciloscopio y guardar la traza completa como archivo \texttt{.csv} (ej. \texttt{medicion\_ambiente.csv}).
\end{itemize}

\begin{resultado}
En la pantalla del osciloscopio, el canal MATH ($CH1 - CH2$) mostrará una senoidal de amplitud reducida con una ligera modulación de ruido de alta frecuencia.\\
Al calentar el termistor NTC con los dedos, la amplitud de la señal en MATH aumentará en tiempo real en la pantalla.\\
El archivo CSV exportado contendrá un encabezado y columnas con los vectores de tiempo y voltaje instantáneo: \texttt{['Time', 'CH1', 'CH2']}.
\end{resultado}

\newpage
% ======================================================================
% FASE 3: PROCESAMIENTO Y VALIDACIÓN EN PYTHON
% ======================================================================
\section*{FASE 3: Procesamiento de Señal y Validación en Python}
El estudiante importa los datos brutos del osciloscopio, filtra ruido, extrae parámetros dinámicos y contrasta los tres mundos (Teoría vs. Simulación vs. Experimento).

\textbf{Paso 3.1 — Ejecución del Script de Validación (\texttt{lab1\_processing.py}):}
\begin{lstlisting}[language=Python]
import numpy as np
import pandas as pd
import matplotlib.pyplot as plt
from scipy.signal import find_peaks

# 1. Carga de datos reales del osciloscopio
df = pd.read_csv("medicion_ambiente.csv")
t = df['Time'].values
v_a = df['CH1'].values
v_b = df['CH2'].values

# Tension diferencial medida
v_out_exp = v_a - v_b

# 2. Extraccion de Amplitud Pico (promediando picos para mitigar ruido)
picos, _ = find_peaks(v_out_exp, height=0.1, distance=50)
v_out_pico_exp = np.mean(v_out_exp[picos])

# 3. Extraccion de Desfase Temporal respecto a Entrada (CH1 representa Vin/2)
picos_in, _ = find_peaks(v_a, height=0.3, distance=50)
delta_t = t[picos[0]] - t[picos_in[0]]
f_c = 10000.0  # 10 kHz
fase_exp_deg = (delta_t * f_c * 360.0) % 360.0

if fase_exp_deg > 180:
    fase_exp_deg -= 360.0

# 4. Modelo Analitico de Referencia (T = 25 C)
v_out_teorico = 0.1815  # V pico
fase_teorica = -162.1   # Grados

# 5. Cuantificacion de Errores
err_amp = np.abs(v_out_pico_exp - v_out_teorico) / v_out_teorico * 100.0
err_fase = np.abs(fase_exp_deg - fase_teorica) / np.abs(fase_teorica) * 100.0

print(f"Amplitud Pico Medida: {v_out_pico_exp:.4f} V | Teorico: {v_out_teorico:.4f} V | Error: {err_amp:.2f}%")
print(f"Fase Angular Medida:  {fase_exp_deg:.2f} deg   | Teorico: {fase_teorica:.2f} deg   | Error: {err_fase:.2f}%")

# 6. Graficacion con Estandar de Publicacion
plt.figure(figsize=(9, 4.5))
plt.plot(t * 1e6, v_out_exp, 'r-', label="Medicion Experimental (Osciloscopio)", alpha=0.8)
t_teo = np.linspace(t[0], t[-1], 1000)
v_teo_plot = v_out_teorico * np.cos(2 * np.pi * f_c * t_teo + np.radians(fase_teorica))
plt.plot(t_teo * 1e6, v_teo_plot, 'k--', label="Modelo Analitico / LTSpice", linewidth=1.5)
plt.xlabel("Tiempo ($\mu s$)")
plt.ylabel("Tension Diferencial $V_{out}$ (V)")
plt.title("Laboratorio 1: Validacion Dinamica del Puente de Wheatstone en CA")
plt.grid(True, linestyle=":")
plt.legend(loc="upper right")
plt.tight_layout()
plt.savefig("lab1_validacion_final.png", dpi=300)
plt.show()
\end{lstlisting}

\begin{resultado}
Una salida en terminal confirmando que el error relativo de amplitud es $\le 5\%$ y el error relativo de fase es $\le 8\%$ (causado principalmente por la tolerancia del capacitor de \qty{4.7}{\nano\farad}).\\
Una figura PNG generada que muestra la curva experimental montada con precisión sobre la curva senoidal teórica calculada.
\end{resultado}

\vspace{0.5cm}
\begin{info}
\textbf{Resumen de la Tabla de Consolidación de Resultados}\\
Cada equipo debe completar la siguiente tabla en su informe técnico final antes de dar por cerrada la práctica.
\end{info}

% --- TABLA ARQUITECTÓNICAMENTE CORREGIDA PARA EVITAR DESBORDAMIENTO ---
\begin{table}[h!]
    \centering
    \small % Reducimos un poco la fuente
    \renewcommand{\arraystretch}{1.5} % Damos aire vertical a las celdas
    \begin{tabularx}{\textwidth}{|>{\raggedright\arraybackslash}p{2.8cm}|>{\centering\arraybackslash}X|>{\centering\arraybackslash}X|>{\centering\arraybackslash}X|>{\centering\arraybackslash}X|>{\centering\arraybackslash}X|}
        \hline
        \textbf{Parámetro Dinámico} & \textbf{Analítico}\newline (Fase 0) & \textbf{LTSpice}\newline (Fase 1) & \textbf{Osciloscopio}\newline (Fase 2) & \textbf{Python}\newline (Fase 3) & \textbf{Error Relativo} \\
        \hline
        $R_{NTC}\ (\qty{25}{\celsius})$ & \qty{10.00}{\kilo\ohm} & \qty{10.00}{\kilo\ohm} & Valor DMM\newline (ej. \qty{9.88}{\kilo\ohm}) & \qty{9.88}{\kilo\ohm} & $< 2.0\%$ \\
        \hline
        $\hat{V}_{out}\ (\qty{25}{\celsius})$ & \qty{181.5}{\milli\volt} & \qty{181.5}{\milli\volt} & $\approx \qty{180}{\milli\volt}$ & \qty{178.2}{\milli\volt} & $< 5.0\%$ \\
        \hline
        Desfase $\theta\ (\qty{25}{\celsius})$ & $-162.1^\circ$ & $-162.1^\circ$ & $\Delta t \approx \qty{45}{\micro\second}$ & $-159.8^\circ$ & $< 6.0\%$ \\
        \hline
        $\hat{V}_{out}\ (\qty{45}{\celsius})$ & \qty{298.3}{\milli\volt} & \qty{298.0}{\milli\volt} & $\approx \qty{305}{\milli\volt}$ & \qty{302.1}{\milli\volt} & $< 5.0\%$ \\
        \hline
    \end{tabularx}
\end{table}

\end{document}